\begin{abstract}
Структурированное состояние обещает ограничить контекст агента, заменив растущую историю взаимодействия компактной памятью. Однако размер промпта не определяет, сколько работы потребуется для решения задачи. Мы исследуем, какие сведения должно сохранять состояние программирующего агента и кто должен их записывать. Сравниваются память, поддерживаемая моделью, гибридное состояние с автоматическими записями среды исполнения и сохранение недавней истории; отдельно проверяются схема и протокол обновления. При замене истории состоянием модели пиковый контекст на общих успехах сокращается до \ModelStatePeakRatio{} от контрольного, но стоимость возрастает в \ModelStateCostRatio{} раза, а общая успешность снижается с \ModelControlRate\% до \ModelStateRate\%. Поисковое сравнение состава гибридной памяти показывает преимущество записей изменений и команд перед одним списком файлов, без устойчивого улучшения от каждого следующего компонента. В повторных измерениях сохранение недавней истории повышает успешность состояния модели до \ModelHistorySixRate\%, гибридной памяти --- до \HybridHistorySixRate\% при \HybridControlRate\% у её собственного контроля. Проверенные изменения схемы и inline-протокол такого восстановления не дают. Стоимость на успех при сохранении истории снижается в первой реализации, но превышает контрольную во второй. Результаты показывают пользу компактной памяти вместе с границами её применения: содержание автоматических записей и доступ к исходным взаимодействиям требуют отдельных решений, а экономию контекста следует оценивать вместе с качеством и стоимостью выполнения.
\end{abstract}
